\vspace{-1em}
\section{Experimental Setup}

\textbf{Dataset Statistics}: We source the claims from QuanTemp and detailed statistics can be found in \cite{quantemp}. It comprises 9935 training, 3084 validation and 2495 test claims. Our evidence collection approach results in 165.7k evidence records.

\textbf{Qualitative Analysis of Decomposed Queries and Evidences }:
To ensure that the queries resulting from \name{} and subsequent evidences collected are of good quality we perform qualitative analysis. Inspired by comprehensiveness measure proposed in \cite{claimdecomp}, we measure the completeness, redundancy and relevance of both the generated queries from \name{} and resulting evidence  with respect to the claim. These measures are defined as follows: 

\textbf{Completeness}: Here the annotators are primarily asked to rate on a scale of (1-5) whether the generated queries cover all aspects of the claim.

\textbf{Redundancy}: The queries generated should be concise and cover independent aspects of the claim (rating 1). if multiple queries cover the same aspect of the claim, they are redundant (rated 5). Depending on extent of redundancy, a rating between 1-5 is assigned. The annotators are also asked to judge the \textbf{Relevance} (0/1) of the generated queries to the claim. Annotators are also asked to rate evidences along the same dimensions. To assess inter-annotator agreement, we calculate Fleiss’ kappa \cite{fleiss}
for each of the annotation tasks. \\
\textbf{Quantitative Analysis of Data Quality}:

Since extensive manual annotation of whole evidence collection is expensive, to automatically evaluate the quality of evidence collected, we compute the upper bound on fact-verification performance attainable using the evidences. We then compare them with  performance attained using gold justification documents. 

\begin{table*}[h]
    \centering
    \begin{tabularx}{\textwidth} { 
       >{\centering\arraybackslash} l 
       c 
       X 
       X 
       X 
       X 
       X
       X 
 }
        \toprule
        \textbf{NLI Classifier} & \textbf{Accuracy} & \multicolumn{3}{c} {\textbf{Per class F1}} & \multicolumn{2}{c}{\textbf{Total}}\\ 
        
         &  & \multicolumn{1}{c}{T} &  \multicolumn{1}{c}{F} & \multicolumn{1}{c}{C} & \multicolumn{1}{c}{W-F1} & \multicolumn{1}{c}{M-F1} \\ \midrule
        \naive & 57.03 & 0.00 & 72.64 & 0.00 & 24.21 & 41.42 \\ 
        \gptt{} & 54.15 & 20.88 & 37.81 & 20.88 & 52.87 & 43.44 \\ 
        \gpttgold{} & 62.32 & 56.67 & 75.35 & 28.00 & 60.47&53.37 \\
        \robertaclaimnoev & 63.2 & 24.71 & 79.93 & 47.37 & 61.53 & 50.67 \\ 
        \robertaqtempplus & \underline{\textbf{\textcolor{gray}{65.25}}} &  \underline{\textbf{\textcolor{gray}{79.92}}} &  \underline{\textbf{\textcolor{gray}{47.76}}} &  \underline{\textbf{\textcolor{gray}{49.67}}} &  \underline{\textbf{\textcolor{gray}{66.56}}} &  \underline{\textbf{\textcolor{gray}{59.11}}} \\ 
        \robertagold & \textbf{69.66} & \textbf{56.86} & \textbf{82.92} & \textbf{48.79} & \textbf{69.79} & \textbf{62.85} \\ \bottomrule
    \end{tabularx}
    \\
    \caption{Quantitative analysis on evidence quality through claim verification performance comparison on \quantempplus.}
    \label{tab:data_quality_quantitative}
    \vspace{-3em}
\end{table*}

    \textbf{ \naive{}}: A Veracity Classifier that only predicts the majority class of the dataset. \robertaclaimnoev fine-tunes a NLI model to predict the veracity of the claim without evidence. 
    
    \textbf{ \gptt{}}: A few-shot learning model that given the claim and corresponding relevant evidence snippets collected using queries from \name{}, predicts the veracity of the claim similar to \cite{quantemp}. We use the \textit{gpt-3.5.-turbo} generative model for this purpose.
    
    \textbf{ \robertaqtempplus{}}: A pre-trained MNLI model fine-tuned to predict the veracity label of the claim given the claim and its corresponding relevant evidence snippets collected using queries from \name{}. \textbf{Please note} for this setting and \gptt{} above we do not perform retrieval from whole collection in open-domain setting as done in retrieval and fact-verification experiments discussed in Sections \ref{5.2} and \ref{5.3}. This setup is  to directly compare quality of evidence collected in \quantempplus{} with gold justification documents from manual fact-checkers.
    
    \textbf{  \robertagold{}}: A pre-trained MNLI model fine-tuned to predict the veracity label of the claim given the claim and its gold justification document authored by human fact-checker. This model serves as the upper bound that can be achieved with the pre-trained MNLI model for \quantempplus{}. 
 \gpttgold{} uses same setups as \robertagold{}, but with \textit{gpt-3.5.-turbo} for veracity prediction.



    \textbf{Fact-Verification:} We employ claim decomposition, followed by retrieval to fetch top-k evidences from our evidence collection in \quantempplus{}. These top-k evidences are then passed on to a NLI model along with the claim to perform verification. We employ Contriever as retrieval model afetr comparing several approaches like BM25, ANCE \cite{ance},Tas-b \cite{tas-b} on validation set.  We experiment with different values of $k=1,3,5,7,10$ for top-k evidence retrieval and observe k=3 to provide best NLI performance on validation set. We fine-tune all the NLI models including ones mentioned above using the Adam optimizer with a weight decay of \( 1\times e^{-5}\) and a learning rate of \(2 \times e^{-5}\). When LLMs are used for verification, we set temperature to 0.1 for reducing randomness in outputs.

    \textbf{Metrics}: For each of these settings, we evaluate the accuracy, per-class F1, macro-F1 (M-F1), and weighted-F1 (W-F1) scores.% to account for the class imbalance in the dataset.

   % \textbf{Hyperparameters}: